% build: latexmk -xelatex cv_zh.tex
% Based on https://github.com/fizixmastr

\documentclass[letterpaper,11pt]{article}
\usepackage[margin=0.5in,bottom=1in,footskip=0.5in]{geometry}
\usepackage{titlesec}
\usepackage{enumitem}
\usepackage[hidelinks]{hyperref}
\usepackage[UTF8]{ctex}


\usepackage{color}
\definecolor{light}{rgb}{0.25, 0.25, 0.25}
\def\light#1{{\color{light}#1}}

\newcommand{\myname}{何思畅}

\usepackage{lastpage}
\usepackage{fancyhdr}
\fancyhf{}
\renewcommand{\headrulewidth}{0pt}
\fancyfoot[L]{\light{\myname{} $\vert{}$ 南加州大学}}
\fancyfoot[R]{\light{第 \thepage{} 页，共 \pageref{LastPage} 页}}
\pagestyle{fancy}


\setlength{\tabcolsep}{0em}

% Sections formatting
\titleformat{\section}{\bfseries\raggedright\large}{}{0em}{}[\titlerule]
\titlespacing*{\section}{0em}{1.1em}{0.25em}
\setlist[itemize]{
  parsep=0em,
  partopsep=0em
}
% • 一级条目，◦ 二级条目
\setlist[itemize,2]{label=\textbullet}
\setlist[itemize,3]{label=\raisebox{0.15ex}{\scriptsize\textopenbullet}}

\newcommand{\CVItem}[1]{
  \item{\small
    {#1}
  }
}

\newcommand{\CVSubheadingFrame}[1]{
  \item
    \begin{tabular*}{0.97\textwidth}[t]{l@{\extracolsep{\fill}}r}
        #1
    \end{tabular*}
}

\newcommand{\CVSubheading}[4]{\CVSubheadingFrame
    {\textbf{#1} & #2 \\
        \small#3 & #4 \\
    }}

\newcommand{\CVSubHeadingListStart}{\begin{itemize}[leftmargin=1.2em, label={}]}
\newcommand{\CVSubHeadingListEnd}{\end{itemize}}
\newcommand{\CVItemListStart}{\begin{itemize}}
\newcommand{\CVItemListEnd}{\end{itemize}}

\definecolor{darkbluepurple}{rgb}{0.125, 0, 0.5}
\newcommand{\uhref}[2]{\href{#1}{\textcolor{darkbluepurple}{#2}}}


% Publication list: [1], [2], ... with hanging indent
\newenvironment{CVRefs}{\begin{list}{}{\setlength{\leftmargin}{2.2em}\setlength{\labelwidth}{1.8em}%
    \setlength{\labelsep}{0.4em}\setlength{\itemsep}{0.3em}\setlength{\parsep}{0em}\usecounter{enumi}%
    \renewcommand{\makelabel}[1]{[\arabic{enumi}]\hfill}}\small}{\end{list}}
\newcommand{\me}{\textbf{何思畅}}
\begin{document}

\noindent
\textbf{\Huge \myname{}}\quad{\large Sichang ``Steven'' He}\hfill
\uhref{mailto:sichangh@usc.edu}
{\small sichangh@usc.edu}\\[0.25em]
计算机科学博士生\hfill
{\small
    \uhref{https://github.com/SichangHe}{GitHub} 与
    \uhref{https://www.youtube.com/@sichanghe}{YouTube}：@SichangHe
}\\
南加州大学网络系统实验室（NSL）\hfill
{\small \uhref{https://sichanghe.github.io/}{sichanghe.github.io}}\\[0.3em]
{\small 关注影响普通人的互联网系统问题。做过大规模互联网测量。}

\section{教育经历}
\CVSubHeadingListStart

\CVSubheadingFrame
{\textbf{南加州大学（USC）}&
    2024年8月 -- 2029年（预计）\quad 美国加州洛杉矶\\
    \small 计算机科学（CS）博士&
    导师：\uhref{https://www.harsha.usc.edu/}{Harsha V. Madhyastha}\\
}
\CVItemListStart
\CVItem
{同时与 \uhref{https://govindan.usc.edu/}{Ramesh Govindan} 合作。}
\CVItem
{部分含项目的课程：
    \begin{itemize}
        \item 高级计算机网络（A+）：
        \uhref{https://github.com/SichangHe/JSphere}{JSphere}：
        测量网页对浏览器 API 的调用，以此对 JavaScript 分类。
    \end{itemize}
    \begin{itemize}
        \item 高级自然语言处理（A）：
            复现了\uhref{https://aclanthology.org/2024.emnlp-main.971/}{一篇
            EMNLP 的大语言模型（LLM）文本检测论文}，并质疑其泛化性。
    \end{itemize}
}
\CVItemListEnd

\CVSubheadingFrame
{\textbf{昆山杜克大学（DKU）与杜克大学双学位}&
    2020年8月 -- 2024年5月\\
    \small 数据科学理学学士（昆山杜克大学授予）与&
    中国昆山与\\
    \small \uhref{https://github.com/SichangHe/sichanghe.github.io/files/15386186/duke_diploma.pdf}
    {跨学科研究理学学士}（方向：数据科学，杜克大学授予）&
    美国北卡达勒姆
}
\CVItemListStart
\CVItem
{修完计算机科学和数学专业的大部分必修课。
}
\CVItem
{部分含项目的课程：
    \begin{itemize}
        \item 计算机网络体系结构（A+）：
        用 Elixir 写的异步 \uhref{https://github.com/SichangHe/CS311}{TCP
            服务器和客户端，以及用 Rust 写的命令行 REPL shell}。

        \item 数据采集与可视化（A）：
        \uhref{https://github.com/SichangHe/STATS401_final_project}{用
            D3.js、Svelte 和 TypeScript 制作的海报}。

        \item 计算机组成与编程（A+，杜克大学）：
        用 MIPS 汇编写二叉搜索树，用 C 写环形缓冲区；
        PizzaCalc 汇编代码长度优化竞赛第 5 名（68 行）。

        \item 数值分析（A+）；抽象代数导论（杜克大学）；
        复分析（A）。
    \end{itemize}
}
\CVItemListEnd

\CVSubheadingFrame
{\textbf{广州市执信中学}&
    2017年8月 -- 2020年6月\quad 中国广州\\
}
\CVItemListStart
\CVItem
{化学奥赛竞赛队；生物奥赛竞赛队。
}
\CVItemListEnd

\CVSubHeadingListEnd

\section{发表论文}
\begin{CVRefs}
    \item \me、Ramesh Govindan、Harsha V. Madhyastha。
    \uhref{https://github.com/SichangHe/DeGenTWeb_docs/releases/tag/cam-ready-imc2025}{Poster: Did I Just Browse A Website Written by LLMs?}
    Proceedings of the 2025 ACM Internet Measurement Conference and Student Workshop (IMC)。
    \uhref{https://www.youtube.com/watch?v=iZHWdpaSOGA}{视频}。2025年10月。
    \item \me、Italo Cunha、Ethan Katz-Bassett。
    \uhref{https://github.com/SichangHe/internet_route_verification/releases/tag/imc-camera-ready}{RPSLyzer: Characterization and Verification of Policies in Internet Routing Registries}。
    Proceedings of the 2024 ACM Internet Measurement Conference (IMC)。
    \uhref{https://youtu.be/xgMPSlQxaH4}{视频}。2024年11月。
    \item \me、Beilong Tang、Boyan Zhang、Jiaqi Shao、Xiaomin Ouyang、
    Daniel Nata Nugraha、Bing Luo。
    \uhref{http://www.arxiv.org/pdf/2402.10464}{FedKit: Enabling Cross-Platform Federated Learning for Android and iOS}。
    IEEE INFOCOM Demo。
    \uhref{https://www.youtube.com/watch?v=TONTBkp_l6M}{视频}。2024年5月。
\end{CVRefs}

\section{在研论文}
\begin{CVRefs}
    \item \me、Calvin Ardi、Ramesh Govindan、Harsha V. Madhyastha。
    \uhref{https://arxiv.org/abs/2605.00087}{DeGenTWeb: A First Look at MGT-dominant Websites}。
    2026年。
\end{CVRefs}

\section{科研经历}
\CVSubHeadingListStart
\CVSubheading {研究助理}{2024年8月 -- 至今}
{南加州大学\uhref{https://nsl.usc.edu/}{网络系统实验室}}
{导师：\uhref{https://www.harsha.usc.edu/}{Harsha V. Madhyastha}}
\CVItemListStart
\CVItem {\textbf{如今网上有多少内容是 LLM 写的？}
    独立完成了首个回答该问题的严谨测量：
    2026 年初存档的网站中，约五分之一主要由 AI 写成。
    用 13.5 万个网站的大规模代表性样本，
    证实了关于网上 LLM 使用情况的坊间说法。
    \begin{itemize}
        \item 构建了 DeGenTWeb：一个准确、可扩展的检测器，
        检测文本内容以 LLM 生成为主的网站。
        \item 在 6 千个基线网站（9 万个网页）上，
        第 90 百分位误报率低至 0.12\%。
        \item 标记出 11.7 万个 Common Crawl 网站中的 21\%，
        以及 Bing 搜索结果中 1.8 万个网站中的 14\%。
        \item 用确凿证据证实了 LLM 文本的使用在增长。
        \item 发现以 LLM 文本为主的网站有明显的盈利动机，
        主要是广告和联盟推广链接。
    \end{itemize}
}
\CVItem {调研了内容来源与真实性联盟（C2PA）标准的使用、弱点和扩展。
}
\CVItemListEnd

\CVSubheading {Rust/Verus 的智能体形式化验证}{2026年6月 -- 2026年10月}
{\uhref{https://www.businesswire.com/news/home/20260205795947/en/CORRECTING-and-REPLACING-Backed-by-OpenAI-Tesla-and-SpaceX-Investors-Midas-Is-Building-Mathematical-Infrastructure-to-Secure-AI-Systems}{Midas
Intelligence Inc.}（面向 AI 的数学初创公司），旧金山（远程）}{}
\CVItemListStart
\CVItem {\textbf{提出了公司的软件路线：}
    借助智能体和 \uhref{https://github.com/verus-lang/verus}{Verus}，
    同时实现开发速度与可靠性。
}
\CVItem {领导并独自创建了公司的软件部门。直接向总裁和 CEO 汇报。
}
\CVItem
{开发了 \uhref{https://github.com/midas-intelligence/midas-lex}{Midas Lex}：
    引导智能体编写 Verus 规约和证明的工具。
}
\CVItem {在 \uhref{https://github.com/uutils/coreutils/issues/13703}{uutils coreutils}
    等热门开源 Rust 项目中发现了 bug。
    构建了未解决 Verus 任务的基准测试集。
}
\CVItemListEnd

\CVSubheading {独立研究：\uhref
    {https://github.com/SichangHe/internet_route_verification}
    {互联网路由验证}
}{2023年4月 -- 2024年5月} {巴西米纳斯吉拉斯联邦大学（远程）}
{导师：\uhref{https://cunha.github.io/about/}{Italo Cunha}}
\CVItemListStart
\CVItem {\textbf{互联网流量是否遵守各网络公开声明的规则？}
    完成了迄今最彻底的检查之一，
    覆盖互联网路由注册库（IRR）中的全部公开路由策略。
    IRR 支撑着互联网的很大一部分：用于指导公共对等互联、
    帮助排查路由问题，并防范可能导致全球断网的路由配置错误和攻击。
}
\CVItem {独立构建了 RPSLyzer：首个能实际使用的解释器，
    可解析 IRR 中全部 7\,GB 路由策略的 99.99\%
    （以路由策略规范语言 RPSL 编写）；
    也是首个按 IETF RFC 2622 和 RFC 4012
    彻底检查真实观测到的边界网关协议（BGP）路由是否遵守这些策略的工具。
    \begin{itemize}
        \item 检查了来自 60 个 RIPE RIS 和 RouteViews 采集点的
        共 7.79 亿条 BGP 路由，耗时 3 小时，内存 2\,GB。
        \item RPSLyzer 能正确解析复杂的 RPSL 结构，
        包括高级集合运算和一种自定义正则表达式。
        \item 刻画了 78,701 个自治系统（AS）的策略。
        发现 RPSL 用法差异很大：有几个 AS 声明了数千条策略，
        而 53\% 一条都没有声明。
    \end{itemize}
}
\CVItem {用 RPSLyzer 检查并分析了观测到的域间路由是否符合路由策略；
    总结了常见的 RPSL 用法模式和用法错误。
    \begin{itemize}
        \item 发现观测到的 BGP 路由中，40\% 的互联因缺少策略而无法验证。
        \item 引入 CAIDA 的 AS 关系数据，使验证覆盖率提高 19 个百分点。
        \item 提供了按 RPSL 验证路由的工具，帮助提升域间路由安全。
        \item 归纳并实现了对 6 种主要原因的检查，
        解释路由为何不符合相关 RPSL。
    \end{itemize}
}
\CVItemListEnd

\CVSubheading {研究助理，移动端联邦学习（FL）}{2023年3月 -- 2024年5月}
{\uhref{https://github.com/FedCampus}{FedCampus 团队}，昆山杜克大学边缘智能实验室}
{导师：\uhref{https://luobing1008.github.io/}{罗冰}}
\CVItemListStart
\CVItem {\textbf{手机能否共同训练 AI，使用设备上的隐私数据，但从不上传这些数据？}
    编写了 \uhref{https://github.com/FedCampus/FedKit}{FedKit SDK}：
    首个可在安卓和 iOS 设备上开展真实联邦学习实验的实用方案，
    让手机无需共享隐私数据即可协作训练共享的机器学习模型。
    \begin{itemize}
    \item 简化研究流程：自动转换模型格式，并把模型直接推送到手机上。
    \end{itemize}
}
\CVItem {编写了 \uhref{https://github.com/flwrlabs/flower/pull/2062}{Flower
    的官方安卓 Kotlin 示例}。Flower 是一个热门联邦学习框架，
    出自剑桥大学研究者创办的 Y Combinator 初创公司。
    因此受邀在 Flower AI Summit 上演讲。
    \begin{itemize}
        \item 发现 iOS 示例\uhref{https://github.com/flwrlabs/flower/issues/2284}{并没有在训练}，
        并协助修正。
    \end{itemize}
}
\CVItem {带领 6 人团队开发了
    \uhref{https://github.com/FedCampus/FedCampus_Flutter}{FedCampus 应用}：
    一款安卓与 iOS 应用，利用 100 多名参与者的个人健康数据，
    在校园里开展真实的联邦学习和联邦分析实验。
    \begin{itemize}
        \item 指导团队成员（4 名本科生和 1 名工程硕士生），其中 2 人由我面试招募。
        \item 调研并主导核心技术选型，包括 Flower、Kotlin 和 Flutter。
    \end{itemize}
}
\CVItemListEnd

\CVSubheading
{研究助理，搜索引擎研究}{2021年12月 -- 2023年5月}
{Search So 团队，昆山杜克大学}{导师：龙江}
\CVItemListStart
\CVItem
{用异步 Rust 开发了一个功能丰富的\uhref{https://github.com/SichangHe/scraper}
    {开源网页爬虫}，用于爬取昆山杜克大学和杜克大学的内网站点。
}
\CVItem
{改进了后端 HTML 处理、前端界面和版本控制。
}
\CVItemListEnd
\CVSubHeadingListEnd

\section{受邀报告}
\begin{CVRefs}
    \item 罗冰、\me。
    \uhref{https://www.youtube.com/watch?v=K7yu2jvu_t8}{FedCampus: a Privacy-preserving Data Platform for Smart Campus with Federated Learning and Analytics}。
    Flower AI Summit，2024年3月。
\end{CVRefs}

\section{教学经历}
\CVSubHeadingListStart
\CVSubheadingFrame
{\textbf{助教，程序设计导论，昆山杜克大学}&
    2021年11月 -- 2022年3月
}
\CVItemListStart
\CVItem
{每周主持实验课；为课程制作了第一个
    \uhref{https://www.youtube.com/watch?v=yiL-ULPBkvE}{开发环境配置视频教程}。
}
\CVItemListEnd

\CVSubheadingFrame
{\textbf{数学与计算机辅导员，昆山杜克大学学术资源中心}&
    2021年5月 -- 2022年5月
}
\CVItemListStart
\CVItem
{获得\uhref
    {https://github.com/SichangHe/curriculum_vitae/files/11665403/CRLA_certificate.pdf}
    {CRLA 国际辅导员培训项目一级认证}。}
\CVItemListEnd
\CVSubHeadingListEnd

\section{学术服务}
\begin{itemize}
    \item 审稿人，IEEE Transactions on Network and Service Management
    （TNSM），2 篇稿件\hfill 2026年
\end{itemize}

\section{奖项}
\begin{itemize}
    \item ACM 互联网测量会议（IMC）差旅资助（2025年9月）
    \item ACM 互联网测量会议（IMC）学生差旅资助（2024年8月）
    \item 通过跑步俱乐部获得昆山杜克大学体育部
    \uhref{https://github.com/SichangHe/curriculum_vitae/files/15110729/scholar_athelete_award_DKU_running_club.pdf}
    {高年级学者运动员奖}（2024年4月）

    \item 国际基因工程机器大赛（iGEM）2022 银牌，昆山杜克大学队（2022年10月）\\{\small
    独自开发了\uhref{https://github.com/SichangHe/igem-2022-dku-backup}
    {团队 wiki}；帮助非技术成员上手 Git；用建模软件验证蛋白质设计。
    }
    \item 昆山杜克大学院长名单（2021年春）及
    优异院长名单（2021年秋、2022年秋、2023年春）
    \item 昆山杜克大学校长奖学金与本科入学奖学金
    （按成绩评定，2020年秋 -- 2024年春）
\end{itemize}

\section{业余项目}
\CVSubHeadingListStart
\CVSubheadingFrame
{\textbf{Rust 库的开源开发者与维护者}&
    2022年11月 -- 至今
}
\CVItemListStart
\CVItem
{\uhref{https://github.com/lzanini/mdbook-katex}{mdBook-KaTeX}：
    \uhref{https://github.com/rust-lang/mdBook}{mdBook} 的数学公式预处理器，
    在 crates.io 上下载量超过 20 万。
    在其无人维护时\uhref{https://github.com/lzanini/mdbook-KaTeX/issues/37}{接手维护}。
    解决了 20 多个历史 issue，新增功能，速度提升 10 倍。
}
\CVItem
{为 \uhref{https://github.com/rust-lang/mdBook}{mdBook}、
    \uhref{https://github.com/cargo-bins/cargo-binstall}{cargo-binstall}、
    \uhref{https://github.com/Byron/pulldown-cmark-to-cmark}{pulldown-cmark-to-cmark}
    等项目贡献代码。
}
\CVItem
{\uhref{https://github.com/SichangHe/agent_manager}{智能体管理器}的%
    \uhref{https://sichangheagentvoicechat.github.io/}{语音控制}、
    \uhref{https://github.com/SichangHe/tokio_gen_server}{Elixir 风格的异步 Rust}、
    \uhref{https://github.com/SichangHe/rust_build_eval}{高效的 Rust 开发环境}、
    \uhref{https://github.com/SichangHe/TeamTypeV}{协作编辑器}。
}
\CVItemListEnd

\CVSubHeadingListEnd

\section{社会服务}
\begin{itemize}
    \item 昆山杜克大学 Intersections 期刊中文编辑与译者\hfill
    2021年6月 -- 2021年8月\\{\small
    合译三篇英文文章为中文；
    审阅编辑多篇文章及
    \uhref{https://sites.duke.edu/intersections/}{Intersections 网站}。
    }
\end{itemize}

\section{技能}
\begin{itemize}[leftmargin=0.5cm, label={}]
    \item \textbf{语言}\quad
    普通话（母语）、粤语（流利）。
    \item \textbf{编程语言}
    \begin{itemize}\small
        \item 深入投入且熟练：Rust、Python、Elixir、JavaScript。
        \item 在项目中用过：
        Kotlin、Dart、Swift、Ruby、Java、Lua、TypeScript、
        Svelte、HTML、CSS、C、SQL、Racket、\LaTeX{}、PL/pgSQL。
        \item 了解：C++、Bash、Julia、Go、Erlang、Gleam、Mathematica、
        MATLAB、Scala、R。
    \end{itemize}
    \item \textbf{并发}\quad
    Rust Tokio、JavaScript、Dart、Swift 中的事件循环；Kotlin 协程；
    Elixir/OTP 进程。
\end{itemize}
\end{document}
